%----------------------------------------------------------------------------------------
% Generic expanded Chinese resume/CV template (final length depends on content)
%----------------------------------------------------------------------------------------
\documentclass[UTF8,a4paper,11pt]{ctexart}

\usepackage{parskip}
\usepackage[margin=0.65in]{geometry}
\usepackage{enumitem}
\usepackage{titlesec}
\usepackage[unicode,hidelinks]{hyperref}

\titleformat{\section}{\Large\bfseries\raggedright}{}{0em}{}[\titlerule]
\titlespacing{\section}{0pt}{6pt}{4pt}
\setlength{\parindent}{0pt}
\hyphenpenalty=5000
\tolerance=1000
\emergencystretch=1em
\clubpenalty=10000
\widowpenalty=10000
\brokenpenalty=10000
\raggedbottom

\newcommand{\cvneedspace}[1]{%
  \par\ifdim\dimexpr\pagegoal-\pagetotal\relax<#1\relax\newpage\fi
}
\newcommand{\cvsection}[1]{\cvneedspace{4\baselineskip}\section{#1}}

\begin{document}
\pagestyle{empty}

\begin{center}
\Huge\textbf{候选人姓名} \\[5pt]
\normalsize
\href{https://github.com/your-handle}{github.com/your-handle} \ $|$ \
\href{mailto:you@example.com}{you@example.com} \ $|$ \
\href{https://www.linkedin.com/in/your-handle}{linkedin.com/in/your-handle}
\end{center}

\cvsection{教育背景}
\cvneedspace{3\baselineskip}
\textbf{学校名称} \hfill \textit{2023年9月 -- 2025年6月}\\
硕士，专业名称

\cvneedspace{3\baselineskip}
\textbf{学校名称} \hfill \textit{2019年9月 -- 2023年6月}\\
学士，专业名称 \hfill 均分 / GPA：按要求选填

\cvsection{经历}
\cvneedspace{4\baselineskip}
\textbf{公司或实验室名称} \hfill \textit{2024年 -- 2025年}\\
\textit{岗位名称}
{\small
\begin{itemize}[leftmargin=1.2em,itemsep=2pt,topsep=3pt]
  \item 用 2 到 4 条说明这个经历解决的核心问题、你的职责边界和交付内容。
  \item 优先写系统能力、模型方案、评测方式、上线边界或业务价值。
  \item 补充有来源支持的结果，例如准确率、时延、覆盖率、测试证据或排名；仅在事实支持时量化。
\end{itemize}
}

\cvneedspace{4\baselineskip}
\textbf{团队或机构名称} \hfill \textit{2023年 -- 2024年}\\
\textit{科研实习生 / 工程师 / 分析师}
{\small
\begin{itemize}[leftmargin=1.2em,itemsep=2pt,topsep=3pt]
  \item 用这一段概括第二段核心经历，强调研究或工程链路中的实际贡献。
\end{itemize}
}

\cvsection{项目经历}
\cvneedspace{4\baselineskip}
\textbf{项目名称} \hfill \textit{2025年}\\
{\small
\begin{itemize}[leftmargin=1.2em,itemsep=2pt,topsep=3pt]
  \item 说明项目目标、核心架构、个人贡献和结果。
  \item 如果是 agent / system / infra 项目，建议明确输入输出、关键模块和安全边界。
\end{itemize}
}

\cvneedspace{4\baselineskip}
\textbf{项目名称} \hfill \textit{2024年}\\
{\small
\begin{itemize}[leftmargin=1.2em,itemsep=2pt,topsep=3pt]
  \item 用更紧凑的方式补充第二个代表项目。
\end{itemize}
}

\cvsection{论文与成果}
{\small
\begin{itemize}[leftmargin=1.2em,itemsep=2pt,topsep=3pt]
  \item \textbf{论文标题}，会议 / 期刊，年份。用一句话说明研究问题、你的关联和价值。
  \item \textbf{成果标题}，比赛 / 数据集 / 基准，年份。用一句话说明结果。
\end{itemize}
}

\cvsection{技能}
\textbf{机器学习 / AI：} 按目标岗位填写可自证技能。\\
\textbf{工程能力：} 按目标岗位填写语言、系统与工具。\\
\textbf{语言 / 其他：} 仅填写真实且相关的能力。

\end{document}
